\section{结果解读：性能爬升、模型差距与计算代价}

\subsection{从低基线到接近人类水平}
讲者展示了通过更好的推理策略和评估机制，系统性能可以从较低水平显著提升，并在部分设置下接近人类表现。这一结果说明 ``系统设计'' 对性能提升同样关键。

\begin{figure}[H]
\centering
\includegraphics[width=0.82\textwidth]{frame-04-human-level.jpg}
\caption{视频约 00:22:23：讲者展示性能提升并提到接近 human level 的区间}
\end{figure}

\begin{importantbox}{这组结果的真正含义}
它并不代表 Agent ``已经解决 Web 自动化''，而是说明在正确的评估和搜索框架下，长链路任务存在明确可优化空间。
\end{importantbox}

\subsection{模型代际演进}
讲者提到 Gemini Pro、GPT-4 以及后续模型在该类任务上的爬升趋势。模型能力的确在进步，但系统中 ``计划-执行-验证'' 的配套仍是必要条件。

\begin{figure}[H]
\centering
\includegraphics[width=0.82\textwidth]{frame-05-model-comparison.jpg}
\caption{视频约 00:25:15：讲者讨论 Gemini Pro、GPT-4 与后续模型在任务表现上的提升}
\end{figure}

\begin{knowledgebox}{如何看待模型提升}
\begin{itemize}
    \item 强模型降低了单步错误率；
    \item 但长链路稳定性仍受搜索和回溯机制影响；
    \item 模型越强，越要配套更严格的评估与治理，避免虚高指标。
\end{itemize}
\end{knowledgebox}

\subsection{成本与收益权衡}
搜索和验证带来成功率提升，同时也增加 token 与环境调用成本。生产系统必须显式管理 ``成功率收益 / 成本开销'' 的平衡。

\begin{warningbox}{成本管理常见失误}
\begin{itemize}
    \item 只追求成功率，忽略单位任务成本；
    \item 忽略失败重试的尾部开销；
    \item 没有任务分级策略，导致简单任务也走重搜索路径。
\end{itemize}
\end{warningbox}

\subsection{本章小结}
实验结果表明 Agent 性能可通过系统化方法显著爬升。但从研究到上线仍需同时优化成功率、时延和成本，而不是单点追分。

